\section*{Hoofdstuk \ref{ch:b1:diagenesis-basins} --- Diagenese en sedimentaire bekkens}

\begin{solution}{exo:b1:diagenesis-basins:1}
Zandsteen $0.49\,e^{-2000/3704} = 0.49\times0.583 = \num{0.286}$; schalie
$0.63\,e^{-2000/1960} = 0.63\times0.360 = \num{0.227}$. De porositeit halveert wanneer
$e^{-z/\lambda} = 1/2$, bij $z = \lambda\ln 2$: \qty{2567}{m} voor zandsteen,
\qty{1359}{m} voor schalie.
\end{solution}

\begin{solution}{exo:b1:diagenesis-basins:2}
$\Delta A = \Delta E + S = 30 + 50 = \qty{80}{m}$. Daarvan wordt \qty{60}{m} met sediment
gevuld, dus $\Delta W = 80 - 60 = \qty{20}{m}$: het plat wordt \qty{20}{m} dieper.
\end{solution}

\begin{solution}{exo:b1:diagenesis-basins:3}
$h = Y(\rho_m - \rho_w)/(\rho_m - \rho_s)$: met \qty{2300}{kg.m^{-3}},
$1.0\times2270/1000 = \qty{2.27}{km}$; met \qty{2500}{kg.m^{-3}},
$1.0\times2270/800 = \qty{2.84}{km}$. Elke meter sediment neemt in de kolom de plaats in van een meter
\omterm{def:b1:layered-earth:layers}{mantel}; hoe dichter het sediment, hoe kleiner de massa die het ten opzichte van de \omterm{def:b1:layered-earth:layers}{mantel} tekort komt,
zodat er meer van nodig is om hetzelfde tekort aan te vullen als het met water gevulde gat.
\end{solution}

\begin{solution}{exo:b1:diagenesis-basins:4}
(a)~Calciet $100.09/2.71 = \qty{36.93}{cm^3.mol^{-1}}$; dolomiet $184.40/2.84
= \qty{64.93}{cm^3.mol^{-1}}$. (b)~Twee mol calciet nemen
\qty{73.86}{cm^3} in beslag, één mol dolomiet \qty{64.93}{cm^3}: de vaste stof krimpt met
$1 - 64.93/73.86 = \qty{12.1}{\%}$. (c)~In een kubieke meter wordt de
\qty{0.90}{m^3} calciet $0.90\times(1 - 0.121) = \qty{0.791}{m^3}$ dolomiet; de
porositeit wordt $1 - 0.791 = \num{0.21}$.
\end{solution}

\begin{solution}{exo:b1:diagenesis-basins:5}
(a)~$s = 1000 - 0.63\times1960\times(1 - e^{-1000/1960}) = \qty{506.5}{m}$; de
gemiddelde porositeit is $1 - 506.5/1000 = \num{0.493}$. (b)~Met de top op
\qty{2000}{m} geeft het oplossen van $s(2000, b) = 506.5$ $b = \qty{2628.7}{m}$: de
laag is \qty{628.7}{m} dik, met gemiddelde porositeit $1 - 506.5/628.7 = \num{0.194}$.
(c)~De korrels zijn dezelfde, dus het verloren water is de verloren dikte:
$1000 - 628.7 = \qty{371}{m^3}$ water per vierkante meter.
\end{solution}

\begin{solution}{exo:b1:diagenesis-basins:6}
(a)~$s = 300 - 0.49\times3704\times(e^{-2500/3704} - e^{-2800/3704}) =
\qty{228.10}{m}$. (b)~Vanaf $b_0 = 300$: $s(0, 300) = 158.80$, $1 - \varphi(300) =
0.5481$, $b_1 = 300 + (228.10 - 158.80)/0.5481 = \qty{426.44}{m}$; dan
$s(0, 426.44) = 229.07$, $1 - \varphi = 0.5633$, $b_2 = 426.44 - 0.96/0.5633 =
\qty{424.73}{m}$; een derde stap verandert hem met minder dan een millimeter. Het
zandsteen was \qty{424.7}{m} dik aan het oppervlak. (c)~$424.7/300 = 1.42$: begraving
heeft het met een factor \num{1.42} dunner gemaakt.
\end{solution}

\begin{solution}{exo:b1:diagenesis-basins:7}
(a)~Het ijs bevat $52\times10^6\times916.7/1000 = \qty{4.77e7}{km^3}$ water;
verdeeld over \qty{361.84e6}{km^2} is het $\qty{0.1317}{km} = \qty{132}{m}$ van het
zeeniveau. (b)~De kustlijnen geven \qty{134}{m}, \qty{2}{m} meer. (c)~$4.767\times
10^7/0.134 = \qty{355.7e6}{km^2}$, \qty{1.7}{\%} minder dan vandaag. Toen het zeeniveau
daalde, kwamen de continentale platen droog te liggen, dus was de glaciale oceaan kleiner, en
verlaagde hetzelfde volume water dat eraan werd onttrokken het zeeniveau meer.
\end{solution}

\begin{solution}{exo:b1:diagenesis-basins:8}
(a)~$(-9.1 - (-98.1))/8 = \qty{11.1}{m}$ per duizend jaar, dicht bij de
\qty{12}{m} uit de kustlijnen. (b)~\qty{40}{mm} per jaar is \qty{40}{m} per
duizend jaar, $40/12 = 3.3$ keer het gemiddelde; in vijfhonderd jaar
$0.040\times500 = \qty{20}{m}$.
\end{solution}

\begin{solution}{exo:b1:diagenesis-basins:9}
(a)~$S = 50\times2 = \qty{100}{m}$ en $\Delta E = \qty{-40}{m}$:
$\Delta A = \qty{60}{m}$, allemaal met sediment gevuld, want de diepte blijft gelijk.
(b)~In \num{20000} jaar zakt het basement slechts $50\times0.02 = \qty{1}{m}$:
het \omterm{def:b1:diagenesis-basins:eustasy}{relatief zeeniveau} daalt met $120 - 1 = \qty{119}{m}$. De zeebodem, die \qty{20}{m} onder het zeeniveau lag, eindigt
$119 - 20 = \qty{99}{m}$ erboven: het hele plat komt droog te liggen.
\end{solution}

\begin{solution}{exo:b1:diagenesis-basins:10}
(a)~Voor een dunne laag is de porositeit er bijna constant: haar vaste
dikte is $s \approx h(1 - \varphi(z))$ op de diepte en $s \approx h_0(1 -
\varphi_0)$ aan het oppervlak. Het gelijkstellen van beide geeft $h_0 \approx h(1 -
\varphi(z))/(1 - \varphi_0)$. (b)~$\varphi(3000) = 0.63\,e^{-3000/1960} =
0.136$, dus $h_0 \approx 10\times0.864/0.37 = \qty{23.3}{m}$, tegen
\qty{23.1}{m} exact. (c)~Tot aan het oppervlak hersteld is de laag
\qty{23}{m} dik, niet langer dun vergeleken met het verval van de porositeit nabij het
oppervlak: haar porositeit daalt van $\varphi_0$ aan haar top tot minder aan haar basis,
dus passen haar korrels in iets minder dan de dikte die overal met $\varphi_0$ is berekend.
\end{solution}

\begin{solution}{exo:b1:diagenesis-basins:11}
(a)~$s = 3000 - 0.63\times1960\times(1 - e^{-3000/1960}) = \qty{2032}{m}$; de
gemiddelde porositeit is $1 - 2032/3000 = \num{0.323}$ en de gemiddelde dichtheid
$0.323\times1030 + 0.677\times2700 = \qty{2161}{kg.m^{-3}}$. (b)~De factor is
$(3300 - 1030)/(3300 - 2161) = \num{1.99}$; de vulling vertegenwoordigt een met water belaste
\omterm{def:b1:diagenesis-basins:subsidence}{subsidentie} van $3000/1.99 = \qty{1505}{m}$. (c)~Met de korreldichtheid zou de factor
$2270/600 = \num{3.78}$ zijn en de \omterm{def:b1:diagenesis-basins:subsidence}{subsidentie} slechts \qty{793}{m}. Dat zou fout zijn omdat
een derde van de vulling water is: het sediment is veel lichter dan
zijn korrels, en alleen de dichtheid van het verzadigde sediment telt in de balans van
massa's.
\end{solution}

\begin{solution}{exo:b1:diagenesis-basins:12}
(a)~De kubus bevat \qty{0.90}{m^3} $= \qty{9.0e5}{cm^3}$ calciet, dat is
$9.0\times10^5/36.93 = \num{24370}$ mol. (b)~Één magnesium per twee calciet:
\num{12180} mol. (c)~Een kilogram zeewater bevat $35.165\times0.0365 =
\qty{1.284}{g}$ magnesium, $1.284/24.305 = \qty{0.0528}{mol}$; de reactie
heeft $12180/0.0528 = \qty{2.31e5}{kg}$ zeewater nodig, $\qty{224}{m^3}$ --- meer
dan tweehonderd keer het volume van het gesteente, zelfs als elk magnesiumion wordt
opgenomen. (d)~Dolomiet kan op grote schaal alleen vormen waar zeer grote volumes
water lang door het gesteente stromen, of waar het water een pekel is die rijker is aan magnesium dan
zeewater: in platforms die door zinkende pekels worden doorgespoeld, bijvoorbeeld, eerder dan in een
sediment dat hetzelfde poriënwater behoudt.
\end{solution}

\begin{solution}{pb:b1:diagenesis-basins:1}
\textbf{1.} Zandsteen: $\varphi(0) = \num{0.490}$, $\varphi(800) =
0.49\,e^{-800/3704} = \num{0.395}$. Schalie: $\varphi(800) = \num{0.419}$,
$\varphi(2600) = \num{0.167}$. Krijt: $\varphi(2600) = \num{0.110}$,
$\varphi(3200) = \num{0.072}$.

\textbf{2.} Met $s$ uit \cref{prop:b1:diagenesis-basins:solid} en $\bar\varphi
= 1 - s/h$ (vraag~7): zandsteen \num{0.441}, schalie \num{0.274}, krijt
\num{0.090}.

\textbf{3.} $h\bar\varphi$: zandsteen \qty{353}{m^3}, schalie \qty{493}{m^3}, krijt
\qty{54}{m^3} water per vierkante meter.

\textbf{4.} $\bar\varphi\rho_w + (1 - \bar\varphi)\rho_g$: zandsteen
$0.441\times1030 + 0.559\times2650 = \qty{1936}{kg.m^{-3}}$, schalie
\qty{2242}{kg.m^{-3}}, krijt \qty{2559}{kg.m^{-3}}. De kolom weegt
$1936\times800 + 2242\times1800 + 2559\times600 = \qty{7.12e6}{kg}$ per vierkante
meter: een gemiddelde dichtheid van \qty{2225}{kg.m^{-3}}.

\textbf{5.} Litostatisch $2225\times9.81\times3200 = \qty{69.9}{MPa}$; hydrostatisch
$1030\times9.81\times3200 = \qty{32.3}{MPa}$; \omterm{def:b1:stress-strain-rheology:lithostatic}{effectieve spanning}
$69.9 - 32.3 = \qty{37.5}{MPa}$.

\textbf{6.} $0.49\,e^{-z/3704} = 0.63\,e^{-z/1960}$ geeft $z = \ln(0.63/0.49)
/(1/1960 - 1/3704) = \qty{1046}{m}$. Beneden die diepte is het zandsteen het poreuzer:
schalie begint natter maar verliest haar water over een kortere afstand.

\textbf{7.} Zandsteen $800 - 0.49\times3704\times(1 - e^{-800/3704}) =
\qty{447.4}{m}$; schalie \qty{1306.7}{m}; krijt \qty{546.0}{m}.

\textbf{8.} \qty{2300.2}{m} korrels; water maakt $1 - 2300.2/3200 =
\qty{28}{\%}$ van de kolom uit.

\textbf{9.} Zandsteen $447.4/15 = \qty{29.8}{m}$, schalie $1306.7/40 =
\qty{32.7}{m}$, krijt $546.0/20 = \qty{27.3}{m}$ vaste stof per miljoen jaar.

\textbf{10.} De schalie had de snelste toevoer, het krijt de langzaamste --- maar de
drie verschillen slechts ongeveer een vijfde, veel minder dan hun dikten vandaag
suggereren.

\textbf{11.} \qty{800}{m}: haar top ligt vandaag aan het oppervlak, waar haar afzetting
net is geëindigd, dus is haar dikte bij afzetting haar dikte vandaag.

\textbf{12.} Het oplossen van $s(0, b) = 1306.7$ voor de schalie met
\cref{met:b1:diagenesis-basins:decompaction} geeft $b = \qty{2124}{m}$: de schalie
was \qty{2124}{m} dik op \qty{15}{\Ma}.

\textbf{13.} $s(0, b) = 546.0$ voor het krijt geeft \qty{1071}{m} op
\qty{55}{\Ma}.

\textbf{14.} Met haar top op \qty{2124}{m} geeft $s(2124, b) = 546.0$ $b =
\qty{2748}{m}$: het krijt was \qty{624}{m} dik.

\textbf{15.} \qty{1071}{m} op \qty{55}{\Ma}; $2124 + 624 = \qty{2748}{m}$ op
\qty{15}{\Ma}.

\textbf{16.} Uit de dikten van vandaag: zandsteen $800/15 = \qty{53.3}{m}$,
schalie $1800/40 = \qty{45.0}{m}$, krijt $600/20 = \qty{30.0}{m}$ per miljoen
jaar. Uit de dikten bij afzetting: \qty{53.3}{m}, $2124/40 =
\qty{53.1}{m}$ en $1071/20 = \qty{53.6}{m}$. De tweede reeks is de zinvolle
--- de snelheid waarmee de zeebodem werd opgebouwd --- en zij toont een constante
snelheid die de dikten van vandaag verbergen.

\textbf{17.} $1 - 1800/2124 = \qty{15}{\%}$.

\textbf{18.} De waterdiepte is aan het begin en aan het eind van elk interval dezelfde, dus is de
gecreëerde \omterm{def:b1:diagenesis-basins:accommodation}{accommodatie} de groei van de kolom boven het referentievlak: \qty{1071}{m} van 75 tot
\qty{55}{\Ma}, $2748 - 1071 = \qty{1677}{m}$ van 55 tot \qty{15}{\Ma}, en $3200 - 2748 = \qty{452}{m}$
sinds \qty{15}{\Ma}.

\textbf{19.} $S = \Delta A - \Delta E$: $1677 - (-60) = \qty{1737}{m}$ van 55 tot
\qty{15}{\Ma}; $452 - (-30) = \qty{482}{m}$ sindsdien.

\textbf{20.} $1737/40 = \qty{43}{m}$ en $482/15 = \qty{32}{m}$ per miljoen jaar.

\textbf{21.} $60/1737 = \qty{3.5}{\%}$ en $30/482 = \qty{6.2}{\%}$. Het zeeniveau daalde in beide intervallen,
dus belemmerde \omterm{def:b1:diagenesis-basins:eustasy}{eustasie} de vulling: het basement moest des te meer zakken om dezelfde ruimte te maken.
\omterm{def:b1:diagenesis-basins:subsidence}{Subsidentie} maakte vrijwel alle \omterm{def:b1:diagenesis-basins:accommodation}{accommodatie}.

\textbf{22.} Schalie en krijt waren \qty{2748}{m} dik op \qty{15}{\Ma} en zijn vandaag \qty{2400}{m}:
hun \omterm{def:b1:diagenesis-basins:compaction}{compactie} maakte \qty{348}{m} ruimte, \qty{77}{\%} van de \qty{452}{m} die door het zandsteen werd gevuld.

\textbf{23.} Elke eenheid is dunner geworden door de begraving die erop volgde, met een hoeveelheid die
afhangt van haar lithologie en van hoe diep zij ging; de ruimte die in elk interval is gemaakt, moet worden
gemeten met de dikten die de eenheden toen hadden, wat alleen \omterm{def:b1:diagenesis-basins:compaction}{decompactie} geeft. En zelfs dan omvat de zo
gevonden \omterm{def:b1:diagenesis-basins:subsidence}{subsidentie} het zinken door het gewicht van het sediment, dat moet worden verwijderd om bij de
\omterm{def:b1:diagenesis-basins:subsidence}{tektonische subsidentie} te komen.

\textbf{24.} \qty{2124}{m}, ongeveer \qty{2.1}{km} --- \qty{324}{m} meer dan de \qty{1800}{m} die de put vandaag vindt.
\end{solution}
